\documentclass[10pt,a4paper]{article}
\usepackage{amsmath,amssymb,amsthm,graphicx}
\usepackage{hyperref}
\usepackage{latexsym}
\usepackage{amssymb}
\usepackage{amsthm}
\usepackage{ulem}
\usepackage{enumitem}
\setlist{itemsep=0.4pt,topsep=0pt}
\parindent 0pt
\parindent 0pt
\parskip 12pt
\newcommand{\HCF}{\mathrm{hcf}}
\newcommand{\ZZ}{\mathbb{Z}}
\newcommand{\rk}{\mathrm{rk}}
\DeclareMathOperator{\len}{length}
\DeclareMathOperator{\bl}{b}
\DeclareMathOperator{\cl}{cl}
\DeclareMathOperator{\di}{d}
\DeclareMathOperator{\Cl}{Cl}
\DeclareMathOperator{\acl}{acl}
\DeclareMathOperator{\tp}{tp}
\DeclareMathOperator{\Th}{Th}
\newcommand{\Cor}{\mathcal{C}_\mathrm{or}}
\newcommand{\Cun}{\mathcal{C}_\mathrm{un}}
\newcommand{\mons}{\mathbb{M}}
\newcommand{\strong}{\mathrm{strong}}
\newcommand{\axioms}{\Delta_{\mathrm{FSTS}}}
\newtheoremstyle{note}% hnamei
{8pt}% hSpace abovei
{14pt}% hSpace belowi
{\normalfont}% hBody fonti
{}% hIndent amounti1
{\bfseries}% hTheorem head fonti
{ }% hPunctuation after theorem headi
{.5em}% hSpace after theorem headi2
{}% hTheorem head spec (can be left empty, meaning ‘normal’)i
\theoremstyle{note}
\newtheorem{thm}{Theorem}[section]
%\newtheorem{theorem}{Theorem}
\newtheorem{lemma}[thm]{Lemma} 
\newtheorem{df}[thm]{Definition}
\newtheorem{rmk}[thm]{Remark} 
\newtheorem{prop}[thm]{Proposition}
\newtheorem{fact}[thm]{Fact}
\newtheorem{cor}[thm]{Corollary}
\newtheorem{conj}[thm]{Conjecture}
\newtheorem{ex}[thm]{Example}
\newtheorem{qu}[thm]{Question}
\newtheorem*{notation}{Notation} 
\title{The closure of a generated subsystem is the subsystem itself}
\author{Silvia Barbina, Enrique Casanovas}
\date{arxiv:2411.13723v3; Creative Commons Attribution 4.0 International}
\begin{document}
\maketitle
\noindent\textit{Editorial note.} The counted statement is the author's labelled environment at source0.tex lines 358-360 and its proof environment opens at line 361; the pairing is by adjacency and the retained proof range is the whole environment.\par\medskip
\noindent\textit{Editorial note.} This source is newly approved for this line; the preamble is built from the author's own package lines, macro block and theorem-environment declarations (source lines 4-65), with the author's optioned amsmath line omitted so that it cannot clash with the wrapper's amsmath.\par\medskip
\noindent\textit{Editorial note.} The source carries no label-number dump and the retained counter values are left empty, so the printed numbering of the card is the wrapper's default; the identity of the counted item is its author label.\par\medskip
\noindent\textit{Editorial note.} Metadata provenance: the source record of the pinned copy and the licence evidence file are both present; the licence is CC BY 4.0 with status supported.\par\medskip
\noindent\textit{Editorial note.} r002 note: the pre-publication gate compares the card's source\_version with the source\_version recorded in the licence evidence of the same source. This source's evidence file records exactly "2411.13723v3", while the card had carried "arXiv:2411.13723v3 (CC BY 4.0)" copied from the source record; r002 aligns the card with the evidence and changes nothing else. The evidence file itself is not modified.\par\medskip
\noindent\textit{Editorial scope.} One original result of the pinned author manuscript, counted once: the closure of a generated subsystem is the subsystem itself (author label order\_closure).\par\medskip
\noindent\textit{Allowed standard prerequisites (editorial).}\begin{itemize}
\item Generated subsystem. The subsystem generated by A is the union of the finite stages obtained by closing under products of two elements already present.. It is the subsystem whose closure is computed.
\item Closure operator. The closure under an ordering is a finitary closure operator and distributes over binary unions.. It is the operator applied to the stages.
\end{itemize}\medskip
\begin{lemma}\label{order_closure}
	If $<$ is an HF-ordering of $M$,  $A\subseteq M$  and $\cl_<(A)=A$, then  $\cl_<(\langle A\rangle )= \langle A \rangle$.
\end{lemma}

\begin{proof}
	Let $\langle A\rangle = \bigcup_{n<\omega}A_n$, where $A_0=A$  and $A_{n+1}$ is the set of elements  of $\langle A\rangle\smallsetminus (A_0\cup \ldots \cup A_n)$  which are a product of two elements of $A_0\cup \ldots 	\cup A_n$.  It is enough to prove by induction that, for every $n$,   $\cl_<(A_0\cup \ldots \cup A_n)= A_0\cup\ldots\cup A_n$. This is clear for $n=0$. Assume $a\in A_{n+1}\smallsetminus (A_0\cup\ldots \cup A_n)$. There are $b,c\in A_0\cup \ldots \cup A_n$  such that $a=b\cdot c$. If $b>a$ or $c>a$, then $a\in \cl_<(A_0\cup \ldots \cup A_n)= A_0\cup \ldots \cup A_n$, a contradiction. Otherwise, $a>b,c$ and $\{b,c\}\subseteq A_0\cup\ldots\cup A_n$ is the unique pair of elements such that $a>b,c$ and $a=b\cdot c$.
\end{proof}

\end{document}
